% !TEX root = ../main.tex
%
% Extension of Oberdieck 2018 (Duke 167:3045) reduced-PT identity
%   Z^{K3 x E}_{red,PT} = -1/Phi_{10}
% to CHL orbifolds (K3 x E)/Z_N, N in {2,3,5,7,11} (Jatkar--Sen index set).
% Audit of the claim that the extension is "seriously open but Pandharipande--
% Thomas technology is actively applicable (not obstructed in principle)".
%
% Scope: five attack-heal cycles; obstacle-type tags (a)/(b)/(c).
% Subscripted \kappa_{ch}, \kappa_{BKM}, \kappa_{cat} throughout.
%
\providecommand{\ClaimStatusTheorem}{\textbf{[Theorem]}}
\providecommand{\ClaimStatusConjectured}{\textbf{[Conjecture]}}
\providecommand{\ClaimStatusAxiomatic}{\textbf{[Axiomatic input]}}
\providecommand{\kappaBKM}{\kappa_{\mathrm{BKM}}}
\providecommand{\kappach}{\kappa_{\mathrm{ch}}}
\providecommand{\kappacat}{\kappa_{\mathrm{cat}}}
\providecommand{\PhiCHL}{\Phi^{\mathrm{CHL}}}
\providecommand{\ZredPT}{Z^{\mathrm{red}}_{\mathrm{PT}}}
\providecommand{\Sp}{\mathrm{Sp}}

\section*{Oberdieck 2018 to CHL-orbifold extension: a reduced-PT roadmap}

Let $X^{(1)} = K3 \times E$ and, for $N \in \{2,3,5,7,11\}$,
\[
  X^{(N)} \;=\; (K3 \times E)/\Z_N,
\]
with $\Z_N$ acting as a balanced symplectic K3-automorphism $g_N$ times
an order-$N$ translation on $E$ (Jatkar--Sen 2005,
arXiv:hep-th/0510147, \S 2; Chaudhuri--Hockney--Lykken 1995). Write
$\ZredPT(X^{(N)})$ for the reduced stable-pairs partition function in
variables $(\tau,\sigma,z)$ conjugate to the K3-fibre class, the
elliptic class, and the point class. The target identity is
\[
  \ZredPT\bigl(X^{(N)}\bigr)(\tau,\sigma,z)
  \;=\; \text{sign} \cdot \frac{1}{\PhiCHL_N(\tau,\sigma,z)},
  \qquad \mathrm{wt}(\PhiCHL_N) = \tfrac{24}{N+1}-2,
\]
with $\PhiCHL_N$ the Jatkar--Sen paramodular form on
$\Sp_4(\Z) \cap \Gamma_0(N)$ (weights $(10,6,4,2,1,0)$ at
$N=1,2,3,5,7,11$). The audit below re-runs the Oberdieck 2018 proof
step by step on the CHL datum and assigns obstacle tags (a) routine
extension, (b) new technique required but achievable, (c) genuine
obstruction pending breakthrough.

\paragraph{Cycle 1: MNOP reduced GW/DT correspondence on $X^{(N)}$.}
MNOP 2006 + Maulik--Pandharipande--Thomas 2010 proves
$Z^{\mathrm{red}}_{\mathrm{GW}} = Z^{\mathrm{red}}_{\mathrm{DT}} =
\ZredPT$ for holomorphic-symplectic CY$_3$. The orbifold extension is
developed by Bryan--Cadman--Young 2012 and Maulik--Pandharipande 2006
(\S 3--4, orbifold MNOP). Orbifold stable pairs exist on
$X^{(N)}$ because $X^{(N)}$ is a Deligne--Mumford stack with trivial
$\omega$ and isolated ADE quotient singularities (the $g_N$ fixed
locus in $K3$ consists of $\chi_g := 24 - \mathrm{fix}(g_N)$ points
contributing type-$A_{N-1}$ cones via the $\Z_N$ action). Tag (a):
the orbifold MNOP correspondence extends routinely by
Bryan--Cadman--Young + Maulik--Pandharipande 2006 Hilbert-scheme
machinery.

\paragraph{Cycle 2: KKV formula on the $g_N$-twisted $K3$.}
The input to $\Phi_{10}$ is the $K3$ elliptic genus
$Z_{\mathrm{ell}}(K3;\tau,z) = 2\phi_{0,1}(\tau,z)$
(Pandharipande--Thomas 2014, Duke 161:2123, stable-pairs proof).
The CHL input is the $g_N$-twined K3 elliptic genus
$\phi^{(g_N)}_{0,1}(\tau,z)$, computed by Cheng--Duncan--Harvey 2014
(arXiv:1307.5793) for all $M_{24}$-classes, with constant term
$c_N(0)$ producing the paramodular weight $c_N(0)/2 = 5,4,3,2,1$ in
the \emph{programme} index set. The Jatkar--Sen weight
$24/(N+1)-2$ differs: at $N=2$, $24/3-2=6$; at $N=3$, $24/4-2=4$. The
twined-KKV proof requires Pandharipande--Thomas 2014 extended to
$g_N$-equivariant stable pairs. Oberdieck--Pandharipande 2019
(Duke 168:1543) constructs equivariant stable pairs for $K3 \times E$
with $E$-translation action and proves equivariant KKV for classes
fixed by the translation. Tag (b): $g_N$-equivariant KKV on $K3$
requires new input beyond PT 2014, but Oberdieck--Pandharipande 2019
\S 5 supplies the equivariant PT machinery; the step is achievable
via $g_N$-equivariant Hilbert-scheme cohomology (Nakajima 1997 +
Lehn 1999 equivariant lift).

\paragraph{Cycle 3: degeneration $E \rightsquigarrow \P^1_{\mathrm{nodal}}$ on the CHL orbifold.}
Oberdieck 2018 \S 3 degenerates $E$ to a nodal $\P^1$ and integrates
the rubber over the node to extract an $\eta^{24}$-factor. On
$X^{(N)}$, the $\Z_N$-translation acts freely on $E$ (no fixed
points) but interacts with the degeneration: the nodal $\P^1$
inherits a $\Z_N$-marked-point structure, and the rubber becomes a
$\Z_N$-equivariant rubber theory. This extends routinely once the
rubber calculus of Maulik--Pandharipande 2006 (\S 3.2,
$\Z_N$-equivariant rubber integrals) is invoked. Tag (a): the
degeneration survives because $E \to E/\Z_N$ is \'etale, so the
nodal-$\P^1$ degeneration lifts without obstruction. The rubber
$\eta^{24}$ factor becomes $\eta^{24/(N+1)}$, reproducing the
Jatkar--Sen weight exponent.

\paragraph{Cycle 4: reduced virtual class on $X^{(N)}$.}
The reduced virtual class on $K3 \times E$ exists because
$H^{0,2}(K3 \times E) = \C$ is one-dimensional and semiregularity
removes it (Kool--Thomas 2014, Oberdieck 2018 \S 2). On $X^{(N)}$,
$H^{0,2}(X^{(N)}) = H^{0,2}(K3)^{g_N} \otimes H^0(E)$; for a
symplectic K3-automorphism the holomorphic 2-form is $g_N$-invariant,
so $H^{0,2}(X^{(N)}) = \C$ persists. The reduced cycle therefore
exists. Tag (a): routine via Kool--Thomas 2014 equivariant cosection
localisation.

\paragraph{Cycle 5: refined/motivic extension $N \geq 2$.}
At $N \geq 2$ the conjectured refined-MNOP identity
$\ZredPT(X^{(N)})^{\mathrm{ref}} = -1/\PhiCHL_N(\text{refined})$
requires a refined GW/DT correspondence. Pandharipande--Pixton--Thomas
2015 (arXiv:1510.03849) gives refined GV on CY$_3$ under the $K3$-fibered assumption; Maulik--Nekrasov 2018 extends to
$\Omega$-deformed instanton counts. The CHL-refined identity has no
Oberdieck-grade theorem but no known obstruction. The hidden
Euler-characteristic consistency:
$\chi_{\mathrm{top}}(K3/g_N) = 24 - \mathrm{fix}(g_N) =
24/(N+1) \cdot (N+1) - \mathrm{fix}(g_N)$; the
$g_N$-twined $\chi$ is $\chi_g = 24/N$ for $N \in \{2,3\}$ (Mukai
1988 Invent. 94), matching the paramodular-weight exponent
$24/(N+1)-2$ after subtracting the two universal factors from the
$(\sigma,\tau)$-Weyl vector. The reduced virtual dimension remains
$1$ on $X^{(N)}$ (Kool--Thomas equivariant). Tag (b) at $N \in \{2,3\}$
(achievable via Oberdieck--Shen 2020, Thomas 2020 $K$-theoretic PT);
tag (c) at $N=11$ where the paramodular weight drops to $0$ and the
cusp form degenerates --- the object $\PhiCHL_{11}$ is a weight-$0$
modular function, not a cusp form, and the reduced PT theory loses
its stable-pairs denominator interpretation.

\paragraph{Summary: consolidated roadmap.}
\ClaimStatusConjectured\ (\emph{CHL-Oberdieck identity}, $N \in \{2,3,5,7\}$)
$\ZredPT(X^{(N)}) = -1/\PhiCHL_N$.
Steps: (a) MNOP orbifold + Kool--Thomas reduced cycle + $E$-degeneration $\Rightarrow$
three of five steps routine. (b) $g_N$-equivariant KKV on $K3$ +
refined CHL at $N \in \{2,3\}$ achievable via
Oberdieck--Pandharipande 2019 \S 5 + Maulik--Nekrasov 2018. (c)
$N=11$ pathology where $\PhiCHL_{11}$ has weight $0$ obstructs
cusp-form interpretation; separate treatment required.
$\kappaBKM(\PhiCHL_N) \ne c_N(0)/2$ for $N \ne 1$ (programme vs.\
CHL dichotomy, Theorem~\ref{thm:cy-c-beyond-k3e-chl-vs-programme-n-list}).
Four subscripted invariants at $X^{(N)}$: $\kappacat(X^{(N)}) = 0$
(Künneth collapses, $E$-factor kills $\chi$);
$\kappach(X^{(N)}) = \sum_q (-1)^q h^{0,q}(X^{(N)})^{g_N}$
equals $\chi(\cO_{X^{(N)}}) = 0$ at $d=3$;
$\kappa_{\mathrm{fiber}}(X^{(N)}) = 24/(N+1)$ (twined K3 Euler).
User's claim that extension is "actively applicable, not obstructed
in principle" is affirmed at $N \in \{2,3,5,7\}$ with obstacle tags
(a)+(b); contradicted at $N=11$ by the weight-zero degeneration
(tag (c)).
